%% system_model.tex
%% Standalone IEEE two-column document: System Model + BBO-DRL Scheduler
%% Target journal: IEEE Journal of Biomedical and Health Informatics (J-BHI)
%% Bio-Inspired Adaptive Task Offloading System for IoT-Edge-Cloud Healthcare
%%
%% Compile with: pdflatex system_model.tex

\documentclass[journal,twoside,letterpaper]{IEEEtran}

%% -----------------------------------------------------------------------
%% Packages
%% -----------------------------------------------------------------------
\usepackage{amsmath}          % align, equation, DeclareMathOperator, etc.
\usepackage{amssymb}          % \mathbb, \mathcal, etc.
\usepackage{bm}               % \bm for bold math
\usepackage{booktabs}         % \toprule, \midrule, \bottomrule
\usepackage{array}            % extended column specifiers in tabular
\usepackage{cite}             % IEEE-style numeric citations
\usepackage{url}              % \url command
\usepackage{graphicx}         % \includegraphics (kept for future figures)
\usepackage{xcolor}           % named colours
\usepackage{hyperref}         % clickable cross-references (load last)
\hypersetup{
  colorlinks = true,
  linkcolor  = blue,
  citecolor  = blue,
  urlcolor   = blue
}

%% -----------------------------------------------------------------------
%% Convenience math operators / macros
%% -----------------------------------------------------------------------
\DeclareMathOperator*{\argmin}{arg\,min}
\newcommand{\vect}[1]{\mathbf{#1}}
\newcommand{\set}[1]{\mathcal{#1}}
\newcommand{\norm}[1]{\left\|#1\right\|}

%% -----------------------------------------------------------------------
%% Title / Author block
%% -----------------------------------------------------------------------
\title{System Model and BBO-DRL Hybrid Scheduler for a\\
       Bio-Inspired Adaptive Task Offloading System\\
       in IoT--Edge--Cloud Healthcare Networks}

\author{%
  \IEEEauthorblockN{Author Name}\\
  \IEEEauthorblockA{Department of Electrical and Computer Engineering\\
    University of Florida, Gainesville, FL 32611, USA\\
    Email: author@ufl.edu}
}

\begin{document}

\maketitle

\begin{abstract}
This document presents the rigorous mathematical system model and the
Bio-inspired Bombardier-Beetle Optimizer -- Deep Reinforcement Learning
(BBO-DRL) hybrid scheduler for adaptive healthcare task offloading in
heterogeneous IoT--Edge--Cloud networks.  The model captures network
topology, task semantics, latency, energy consumption, privacy risk, and
patient criticality, and culminates in a constrained multi-objective
optimisation problem.  The BBO-DRL framework solves this problem online
with provable real-time complexity.  Full derivations, parameter
definitions, and hardware profiles are provided.
\end{abstract}

\begin{IEEEkeywords}
task offloading, mobile edge computing, deep reinforcement learning,
bio-inspired optimisation, IoMT, healthcare, privacy, patient criticality.
\end{IEEEkeywords}

%% =======================================================================
\section{System Model}
\label{sec:system_model}
%% =======================================================================

\subsection{Network and Task Definitions}
\label{subsec:network_task}

We consider a three-tier heterogeneous network composed of Internet-of-Medical-Things
(IoMT) wearable devices, edge/fog infrastructure, and a remote cloud server.

\subsubsection{Network Topology}

Let
\begin{align}
  \set{U} &= \{u_1, u_2, \ldots, u_N\}, \label{eq:set_U}\\
  \set{E} &= \{e_1, e_2, \ldots, e_M\}, \label{eq:set_E}\\
  \set{C} &= \{c\},                      \label{eq:set_C}
\end{align}
denote the set of $N$ IoMT wearable devices, the set of $M$ edge/fog compute
nodes (edge gateways and fog servers), and the single cloud server,
respectively.  All devices in $\set{U}$ communicate over a shared wireless
channel characterised by bandwidth $B$ (Hz), while edge/fog nodes are
interconnected via high-throughput wired or 5G backhaul links.

\subsubsection{Healthcare Task Model}

A healthcare task generated by device $u_i \in \set{U}$ at time $t_i$ is
described by the 4-tuple
\begin{equation}
  \tau_i = \left(D_i,\; C_i,\; T^{\max}_i,\; \rho_i\right),
  \label{eq:task_tuple}
\end{equation}
where
\begin{itemize}
  \item $D_i$ (bits) is the data payload to be transmitted for remote execution;
  \item $C_i$ (CPU cycles) is the total number of processor cycles required to
        complete the task;
  \item $T^{\max}_i$ (seconds) is the maximum tolerable end-to-end delay,
        defined by the clinical service-level agreement (SLA) for task type
        $\tau_i$ (e.g., 150\,ms for ECG anomaly detection, 500\,ms for
        activity recognition);
  \item $\rho_i \in [0,1]$ is the privacy sensitivity level of the health data
        payload, with $\rho_i = 1$ indicating maximally sensitive information
        (e.g., full EHR records) and $\rho_i = 0$ indicating non-sensitive
        telemetry.
\end{itemize}

\subsubsection{Offloading Decision Variable}

For each task $\tau_i$ the scheduler selects an execution node encoded by
the integer decision variable
\begin{equation}
  x_i \;\in\; \{0,\, 1,\, 2,\, \ldots,\, M+2\},
  \label{eq:decision_var}
\end{equation}
where $x_i = 0$ denotes local execution on the wearable $u_i$,
$x_i = 1$ selects the nearest edge gateway, $x_i \in \{2, \ldots, M+1\}$
selects fog node $e_{x_i - 1}$, and $x_i = M+2$ routes the task to the
cloud server $c$.

%% -----------------------------------------------------------------------
\subsection{Latency Model}
\label{subsec:latency}
%% -----------------------------------------------------------------------

\subsubsection{Local Execution Delay}

When $x_i = 0$, task $\tau_i$ executes entirely on the wearable.  The
local computation delay is
\begin{equation}
  T^{\mathrm{local}}_i = \frac{C_i}{f_w},
  \label{eq:latency_local}
\end{equation}
where $f_w$ (cycles/s) is the CPU clock frequency of the wearable device.

\subsubsection{Wireless Uplink Data Rate}

For offloading to node $j$ ($j \neq 0$), device $u_i$ transmits over a
wireless link.  Applying Shannon's channel-capacity theorem, the achievable
uplink rate is
\begin{equation}
  R_{i,j} = B \log_2\!\left(1 + \frac{P^t_i \, h_{i,j}}{\sigma^2 + I_j}\right),
  \label{eq:shannon_rate}
\end{equation}
where $P^t_i$ (W) is the transmission power of $u_i$, $h_{i,j}$ is the
channel power gain between $u_i$ and node $j$, $\sigma^2$ (W) is the
additive white Gaussian noise (AWGN) power spectral density integrated over
$B$, and $I_j$ (W) is the aggregate co-channel interference received at
node $j$ from all other simultaneously transmitting devices.

\subsubsection{Path-Loss Channel Gain}

The channel gain $h_{i,j}$ is modelled using a log-distance path-loss law:
\begin{equation}
  h_{i,j} = h_0 \left(\frac{d_0}{d_{i,j}}\right)^{\!\alpha},
  \label{eq:channel_gain}
\end{equation}
where $d_{i,j}$ (m) is the Euclidean distance between $u_i$ and node $j$,
$d_0$ (m) is a near-field reference distance, $h_0$ is the reference channel
gain measured at $d_0$, and $\alpha \in [2, 4]$ is the environment-dependent
path-loss exponent ($\alpha = 2$ for free-space, $\alpha = 3$--$4$ for indoor
clinical environments).

\subsubsection{Queuing Delay at Compute Nodes}

Task arrivals at each compute node are modelled as a Poisson process with
mean arrival rate $\lambda_j$ (tasks/s) and exponentially distributed service
times with mean rate $\mu_j$ (tasks/s), forming an M/M/1 queue.  The
expected waiting time in queue is given by the Pollaczek--Khinchine formula:
\begin{equation}
  T^{\mathrm{queue}}_{i,j} = \frac{\lambda_j}{\mu_j \left(\mu_j - \lambda_j\right)},
  \quad \lambda_j < \mu_j,
  \label{eq:queue_delay}
\end{equation}
where the stability condition $\lambda_j < \mu_j$ must hold; tasks submitted
to a saturated node ($\lambda_j \geq \mu_j$) are rejected by the admission
controller and reassigned by the scheduler.

\subsubsection{Total Offloading Latency}

The end-to-end offloading latency from device $u_i$ to compute node $j$
comprises the wireless transmission delay, the electromagnetic propagation
delay, the queueing wait, and the remote computation time:
\begin{equation}
  T^{\mathrm{offload}}_{i,j} =
    \underbrace{\frac{D_i}{R_{i,j}}}_{\text{transmission}}
    + \underbrace{\frac{d_{i,j}}{v_s}}_{\text{propagation}}
    + \underbrace{T^{\mathrm{queue}}_{i,j}}_{\text{queueing}}
    + \underbrace{\frac{C_i}{f_j}}_{\text{computation}},
  \label{eq:latency_offload}
\end{equation}
where $v_s \approx 3 \times 10^8$\,m/s is the electromagnetic signal
propagation speed and $f_j$ (cycles/s) is the clock frequency of compute
node $j$.  Result return latency is omitted from the SLA-critical path
because the gateway pre-fetches results asynchronously; however, it is
included in the energy model (Section~\ref{subsec:energy}).

\subsubsection{End-to-End Task Latency}

Combining local and offloading cases, the end-to-end latency for task
$\tau_i$ under decision $x_i$ is
\begin{equation}
  T_i(x_i) =
  \begin{cases}
    T^{\mathrm{local}}_i                   & \text{if } x_i = 0, \\[4pt]
    T^{\mathrm{offload}}_{i,\, x_i}       & \text{if } x_i \geq 1.
  \end{cases}
  \label{eq:latency_total}
\end{equation}

%% -----------------------------------------------------------------------
\subsection{Energy Consumption Model}
\label{subsec:energy}
%% -----------------------------------------------------------------------

Energy expenditure is modelled from the perspective of the resource-constrained
wearable $u_i$, as the wearable's battery lifetime is the primary hardware
constraint in continuous patient monitoring.

\subsubsection{Local Computation Energy (CMOS Model)}

For local execution the dynamic power dissipation follows the standard CMOS
switched-capacitance model.  The total energy consumed is
\begin{equation}
  E^{\mathrm{local}}_i = \kappa \, C_i \, f_w^2,
  \label{eq:energy_local}
\end{equation}
where $\kappa$ (F/cycle$^2$, or equivalently J\,s$^2$/cycle$^3$) is the
effective switched capacitance, a hardware architecture constant that
encapsulates the average activity factor and load capacitance per clock
cycle~\cite{burd1995processor}.  Representative values are listed in
Table~\ref{tab:hardware}.

\subsubsection{Transmission and Reception Times}

The time required for $u_i$ to upload the task payload to node $j$ is
\begin{equation}
  T^{\mathrm{tx}}_{i,j} = \frac{D_i}{R_{i,j}},
  \label{eq:time_tx}
\end{equation}
and the time to receive the compressed result of size $D^{\mathrm{ret}}_i$
(bits) back from node $j$ over the downlink (assumed symmetric) is
\begin{equation}
  T^{\mathrm{rx}}_{i,j} = \frac{D^{\mathrm{ret}}_i}{R_{i,j}}.
  \label{eq:time_rx}
\end{equation}

\subsubsection{Wearable Offloading Energy}

During offloading the wearable enters two sub-states: \emph{active
transmission} (duration $T^{\mathrm{tx}}_{i,j}$) and \emph{idle listening}
(duration $T^{\mathrm{offload}}_{i,j} - T^{\mathrm{tx}}_{i,j}$, covering
propagation, queueing, remote computation, and result download).  The total
wearable energy is
\begin{equation}
  E^{\mathrm{offload}}_i =
    P^t_i \cdot T^{\mathrm{tx}}_{i,j}
    + P_{\mathrm{idle}} \cdot \!\left(T^{\mathrm{offload}}_{i,j}
      - T^{\mathrm{tx}}_{i,j}\right),
  \label{eq:energy_offload}
\end{equation}
where $P_{\mathrm{idle}}$ (W) is the quiescent (listening/sleep) power draw
of the wearable's radio and CPU subsystem.

\subsubsection{Total Wearable Energy}

\begin{equation}
  E_i(x_i) =
  \begin{cases}
    E^{\mathrm{local}}_i   & \text{if } x_i = 0, \\[4pt]
    E^{\mathrm{offload}}_i & \text{if } x_i \geq 1.
  \end{cases}
  \label{eq:energy_total}
\end{equation}

%% -----------------------------------------------------------------------
\subsection{Privacy Risk Quantification}
\label{subsec:privacy}
%% -----------------------------------------------------------------------

Repeated offloading of health data to a fixed node creates a predictable
routing pattern that may be exploited by a passive traffic-analysis attacker
to infer patient identity or clinical state.  We quantify this risk using
information-theoretic entropy.

\subsubsection{Offloading Probability Distribution}

Let $p_{i,j}$ denote the empirical probability that device $u_i$ offloads a
task to compute node $e_j$, estimated over a sliding observation window of
the most recent $W$ scheduling decisions.  By construction,
\begin{equation}
  \sum_{j=0}^{M+2} p_{i,j} = 1, \quad p_{i,j} \geq 0 \;\; \forall\, j.
  \label{eq:prob_dist}
\end{equation}

\subsubsection{Shannon Routing Entropy}

The Shannon entropy of the routing distribution for device $u_i$ is
\begin{equation}
  H(u_i) = -\sum_{\substack{j=0 \\ p_{i,j} > 0}}^{M+2}
              p_{i,j} \log_2 p_{i,j} \quad \text{(bits)},
  \label{eq:shannon_entropy}
\end{equation}
which measures the uncertainty available to an attacker attempting to predict
the next offloading destination.  High entropy corresponds to an unpredictable
routing pattern that is difficult to fingerprint.

\subsubsection{Maximum Entropy Baseline}

The upper bound is achieved when the task load is perfectly balanced across
all $|\set{E}|+1$ possible destinations (including local):
\begin{equation}
  H_{\max} = \log_2\!\left(|\set{E}| + 1\right) \quad \text{(bits)}.
  \label{eq:entropy_max}
\end{equation}

\subsubsection{Privacy Risk Cost}

We define the per-task privacy risk cost as
\begin{equation}
  R_P(u_i, x_i) = \rho_i \cdot \left(1 - \frac{H(u_i)}{H_{\max}}\right),
  \label{eq:privacy_risk}
\end{equation}
where $\rho_i$ is the data sensitivity level from~\eqref{eq:task_tuple}.
Equation~\eqref{eq:privacy_risk} satisfies the following boundary conditions:
\begin{itemize}
  \item $R_P = 0$ when the routing entropy is maximal ($H(u_i) = H_{\max}$),
        i.e., load is perfectly distributed and the attacker gains no
        information advantage;
  \item $R_P = \rho_i$ when all tasks are routed to the same node
        ($H(u_i) = 0$), i.e., the wearable's data path is fully predictable.
\end{itemize}

%% -----------------------------------------------------------------------
\subsection{Patient Criticality Index}
\label{subsec:criticality}
%% -----------------------------------------------------------------------

\subsubsection{Criticality Index Definition}

A dedicated Criticality Index (CI) module processes the raw physiological
streams (ECG, SpO$_2$, blood pressure, accelerometer) and computes a
normalised scalar
\begin{equation}
  \Phi_i \in [0, 1],
  \label{eq:ci_def}
\end{equation}
where $\Phi_i = 0$ denotes a healthy baseline and $\Phi_i = 1$ indicates
a life-threatening emergency.  The CI module employs a fuzzy-rule physiological
anomaly scoring engine; its internal design is outside the scope of this
section.

\subsubsection{CI-Adaptive Weight Functions}

The multi-objective cost in Section~\ref{subsec:cost} balances three competing
objectives: energy efficiency, latency minimisation, and privacy preservation.
As patient criticality rises, clinical urgency demands that latency be
prioritised over energy conservation and privacy.  We model this adaptive
trade-off through three continuous, non-linear weight functions of $\Phi_i$:

\paragraph{Energy weight (exponential decay)}
\begin{equation}
  w_E(\Phi_i) = e^{-\alpha_E \Phi_i},
  \label{eq:weight_energy}
\end{equation}
which decreases monotonically from $1$ (at $\Phi_i = 0$) towards $0$ as
criticality increases, reflecting that battery conservation becomes secondary
during emergencies.

\paragraph{Latency weight (normalised exponential growth)}
\begin{equation}
  w_L(\Phi_i) = \frac{e^{\beta_L \Phi_i} - 1}{e^{\beta_L} - 1},
  \label{eq:weight_latency}
\end{equation}
which grows from $0$ to $1$ over $[0,1]$ and is normalised so that
$w_L(1) = 1$ exactly, ensuring that a critically ill patient always
receives maximum latency priority.

\paragraph{Privacy weight (polynomial decay)}
\begin{equation}
  w_P(\Phi_i) = (1 - \Phi_i)^{\gamma_P},
  \label{eq:weight_privacy}
\end{equation}
which decays polynomially from $1$ to $0$, modelling the clinical ethic that
privacy constraints may be relaxed (e.g., emergency data may be routed to
the nearest available node) when patient survival is at stake.

\paragraph{Normalised weights}

To guarantee that the three weights sum to unity for any $\Phi_i$, we define
the normalisation factor
\begin{equation}
  W(\Phi_i) = w_E(\Phi_i) + w_L(\Phi_i) + w_P(\Phi_i),
  \label{eq:weight_sum}
\end{equation}
and compute the final scheduler weights as
\begin{equation}
  \hat{w}_k(\Phi_i) = \frac{w_k(\Phi_i)}{W(\Phi_i)},
  \quad k \in \{E,\, L,\, P\},
  \label{eq:weight_normalised}
\end{equation}
satisfying $\hat{w}_E + \hat{w}_L + \hat{w}_P = 1$ by construction.

\paragraph{Hyperparameter values}

The shape parameters $\alpha_E = 3$, $\beta_L = 4$, and $\gamma_P = 2$ are
set empirically via grid search on a held-out clinical trace dataset.  The
sensitivity of system performance to these parameters is analysed in
Section~V of the full paper; results show $<\!3\%$ variation in normalised
cost across $\pm 30\%$ perturbations, confirming robustness.

%% -----------------------------------------------------------------------
\subsection{Multi-Objective Cost Function and Optimisation Problem}
\label{subsec:cost}
%% -----------------------------------------------------------------------

\subsubsection{Normalised Performance Metrics}

To render the three objectives commensurable, each raw metric is normalised
via min-max scaling over the finite set of candidate execution nodes
$\set{X} = \{0, 1, \ldots, M+2\}$:
\begin{align}
  \hat{E}_i(x) &= \frac{E_i(x) - E^{\min}}{E^{\max} - E^{\min}},
  \label{eq:norm_energy}\\[4pt]
  \hat{T}_i(x) &= \frac{T_i(x) - T^{\min}}{T^{\max}_{\mathrm{range}} - T^{\min}},
  \label{eq:norm_latency}\\[4pt]
  \hat{R}_P(x) &= R_P(u_i, x),
  \label{eq:norm_privacy}
\end{align}
where $E^{\min} = \min_{x \in \set{X}} E_i(x)$,
$E^{\max} = \max_{x \in \set{X}} E_i(x)$, and similarly for
$T^{\min}$, $T^{\max}_{\mathrm{range}}$.  Note that $R_P \in [0,1]$
already (see~\eqref{eq:privacy_risk}), so no further scaling is needed.

\subsubsection{Aggregate Cost Function}

The scalar cost to be minimised by the scheduler is the CI-weighted linear
combination of the three normalised objectives:
\begin{equation}
  F(x_i) =
    \hat{w}_E(\Phi_i) \cdot \hat{E}_i(x_i)
    + \hat{w}_L(\Phi_i) \cdot \hat{T}_i(x_i)
    + \hat{w}_P(\Phi_i) \cdot \hat{R}_P(x_i).
  \label{eq:cost_function}
\end{equation}
Because all three terms lie in $[0,1]$ and the weights are normalised,
$F(x_i) \in [0,1]$.

\subsubsection{Constrained Optimisation Problem}

The per-task scheduling problem is formulated as
\begin{align}
  &\minimize_{x_i}\; F(x_i) \label{eq:opt_obj}\\[4pt]
  &\text{subject to:}\nonumber\\
  &\quad T_i(x_i) \leq T^{\max}_i,     \label{eq:constraint_sla}\\
  &\quad E_i(x_i) \leq E^{\mathrm{budget}}_i, \label{eq:constraint_energy}\\
  &\quad x_i \in \{0, 1, \ldots, M+2\}, \label{eq:constraint_domain}
\end{align}
where constraint~\eqref{eq:constraint_sla} enforces the clinical SLA,
constraint~\eqref{eq:constraint_energy} ensures the wearable's instantaneous
energy draw does not exhaust the battery budget $E^{\mathrm{budget}}_i$
allocated for this sensing cycle, and~\eqref{eq:constraint_domain} restricts
the decision to valid integer node indices.

\medskip
Table~\ref{tab:hardware} summarises the hardware parameters used to
instantiate~\eqref{eq:latency_local}--\eqref{eq:constraint_domain} in
simulation.

%% -----------------------------------------------------------------------
%% Hardware Parameters Table
%% -----------------------------------------------------------------------
\begin{table*}[t]
  \centering
  \caption{Hardware Platform Parameters Used in Simulation}
  \label{tab:hardware}
  \renewcommand{\arraystretch}{1.25}
  \begin{tabular}{%
      l
      >{\centering\arraybackslash}p{2.8cm}
      >{\centering\arraybackslash}p{2.8cm}
      >{\centering\arraybackslash}p{2.5cm}
      >{\centering\arraybackslash}p{2.5cm}}
    \toprule
    \textbf{Parameter} &
    \textbf{Wearable (ESP32-S3)} &
    \textbf{Edge Gateway (RPi\,4)} &
    \textbf{Fog Node} &
    \textbf{Cloud Server} \\
    \midrule
    CPU frequency $f$ (MHz)
      & 240 & 1\,500 & 2\,200 & 3\,200 \\
    Eff.\ capacitance $\kappa$ (F/cycle$^2$)
      & $10^{-27}$ & $10^{-28}$ & $10^{-28}$ & --- \\
    Tx power $P^t$ (mW)
      & 178 & --- & --- & --- \\
    Idle power $P_{\mathrm{idle}}$ (mW)
      & 33 & 3\,400 & 5\,000 & --- \\
    RAM
      & 512\,KB & 8\,GB & 16\,GB & $\infty$ \\
    Channel bandwidth $B$
      & 20\,MHz (Wi-Fi) & 100\,MHz (5G) & 100\,MHz & 1\,Gbps \\
    \bottomrule
  \end{tabular}
  \smallskip\\
  {\footnotesize ``---'' indicates the parameter is not applicable or
  not a binding constraint for that tier.  Cloud RAM is treated as
  effectively unlimited at the task-level granularity considered here.}
\end{table*}

%% =======================================================================
\section{The BBO-DRL Hybrid Scheduler}
\label{sec:bbo_drl}
%% =======================================================================

\subsection{Algorithm Overview}
\label{subsec:algo_overview}

The constrained integer optimisation problem in~\eqref{eq:opt_obj}--\eqref{eq:constraint_domain}
is NP-hard in general (reducible from the multi-dimensional knapsack problem).
We propose a two-phase hybrid algorithm that combines:
\begin{enumerate}
  \item \textbf{Bombardier Beetle Optimizer (BBO)}: a population-based
        meta-heuristic that performs combinatorial search over offloading
        routes, inspired by the \emph{Brachinus} beetle's ballistic chemical
        spray and predator-evasion behaviour; and
  \item \textbf{Deep Q-Network (DQN)}: a model-free deep reinforcement
        learning agent that encodes real-time environment state (channel
        quality, compute load, battery level, patient criticality) to guide
        the BBO search towards a dynamically relevant subspace, dramatically
        reducing the effective search dimensionality.
\end{enumerate}
The DQN pre-filters candidate nodes; the BBO exploits that reduced space
to find a near-optimal $x_i^*$ within the latency budget of a single
scheduling epoch.

%% -----------------------------------------------------------------------
\subsection{Bombardier Beetle Optimizer (BBO) Formulation}
\label{subsec:bbo}
%% -----------------------------------------------------------------------

\subsubsection{Population Representation}

BBO maintains a population of $N_{\mathrm{pop}}$ candidate solutions (beetles),
where the $k$-th beetle is represented as a real-valued parameter vector
$\vect{X}_k \in \mathbb{R}^d$.  Each dimension encodes a continuous
relaxation of an offloading sub-decision; the final integer $x_i$ is
recovered by projection onto $\{0, \ldots, M+2\}$ after optimisation.

\subsubsection{Exploration Phase: Chemical Spray}

Inspired by the beetle's long-range explosive chemical jet, the exploration
update generates a large stochastic perturbation anchored to the current
global best $\vect{X}_{\mathrm{best}}(t)$:
\begin{equation}
  \vect{X}^{\mathrm{spray}}_k(t+1) =
    \vect{X}_k(t)
    + \vect{r}_1 \odot \!\left(\vect{X}_{\mathrm{best}}(t) - \vect{X}_k(t)\right)
    + r_2 \cdot \delta(t) \cdot \mathrm{sign}\!\left(\vect{u} - 0.5\right),
  \label{eq:bbo_spray}
\end{equation}
where $\vect{r}_1, \vect{u} \sim \mathcal{U}(0,1)^d$ are independent
uniform random vectors, $r_2 \sim \mathcal{U}(0,1)$ is a scalar random
variable, $\odot$ denotes element-wise (Hadamard) multiplication, and
$\delta(t)$ is the iteration-dependent spray radius defined
in~\eqref{eq:delta_decay} below.  The $\mathrm{sign}(\cdot)$ term
introduces a symmetric random directional bias that promotes diverse
exploration of the feasible space.

\subsubsection{Exploitation Phase: Predator Avoidance}

The beetle's close-range evasion manoeuvre drives fine-grained local search
towards $\vect{X}_{\mathrm{best}}(t)$:
\begin{equation}
  \vect{X}^{\mathrm{avoid}}_k(t+1) =
    \vect{X}_k(t)
    + r_3 \cdot \mathrm{step}_k(t) \cdot
      \frac{\vect{X}_{\mathrm{best}}(t) - \vect{X}_k(t)}
           {\norm{\vect{X}_{\mathrm{best}}(t) - \vect{X}_k(t)} + \varepsilon},
  \label{eq:bbo_avoid}
\end{equation}
where $r_3 \sim \mathcal{U}(0,1)$, $\varepsilon > 0$ prevents division by
zero, and $\mathrm{step}_k(t)$ is a fitness-proportionate step size:
\begin{equation}
  \mathrm{step}_k(t) =
    \frac{\norm{\vect{X}_k(t) - \vect{X}_{\mathrm{predator}}}}
         {\norm{\vect{X}_{\mathrm{best}}(t) - \vect{X}_{\mathrm{predator}}} + \varepsilon}.
  \label{eq:bbo_step}
\end{equation}
Here $\vect{X}_{\mathrm{predator}}$ is a stochastically placed adversarial
reference point that diversifies the local search direction across the
population.

\subsubsection{Greedy Selection}

After computing both candidate updates, the beetle retains the one with
lower cost:
\begin{equation}
  \vect{X}_k(t+1) =
  \begin{cases}
    \vect{X}^{\mathrm{spray}}_k(t+1)
      & \text{if } F\!\left(\vect{X}^{\mathrm{spray}}_k(t+1)\right)
                 < F\!\left(\vect{X}^{\mathrm{avoid}}_k(t+1)\right),\\[4pt]
    \vect{X}^{\mathrm{avoid}}_k(t+1) & \text{otherwise.}
  \end{cases}
  \label{eq:bbo_selection}
\end{equation}

\subsubsection{Spray Radius Decay}

To transition from global exploration early in the search to fine
exploitation near convergence, the spray radius follows a quadratic decay
schedule:
\begin{equation}
  \delta(t) = \delta_0 \left(1 - \frac{t}{T_{\mathrm{BBO}}}\right)^{\!2},
  \label{eq:delta_decay}
\end{equation}
where $\delta_0$ is the initial (maximum) radius and $T_{\mathrm{BBO}}$ is
the total number of BBO iterations per scheduling epoch.  At $t = 0$,
$\delta = \delta_0$ (maximum diversity); at $t = T_{\mathrm{BBO}}$,
$\delta = 0$ (pure exploitation).

%% -----------------------------------------------------------------------
\subsection{Deep Q-Network (DQN) State-Action Framework}
\label{subsec:dqn}
%% -----------------------------------------------------------------------

\subsubsection{State Space}

The DQN agent observes a six-dimensional state vector at each scheduling
epoch $t$:
\begin{equation}
  \vect{s}_t = \left[
    \mathrm{BW}_t,\;
    \mathrm{CPU}_t,\;
    \mathrm{RTT}_t,\;
    \mathrm{bat}_t,\;
    \Phi_t,\;
    p^{\mathrm{atk}}_t
  \right]^\top \in [0,1]^6,
  \label{eq:dqn_state}
\end{equation}
where all components are min-max normalised to $[0,1]$:
$\mathrm{BW}_t$ = available channel bandwidth utilisation,
$\mathrm{CPU}_t$ = average CPU utilisation across edge/fog nodes,
$\mathrm{RTT}_t$ = measured round-trip time to the best available node,
$\mathrm{bat}_t$ = remaining wearable battery fraction,
$\Phi_t$ = current patient criticality index from~\eqref{eq:ci_def}, and
$p^{\mathrm{atk}}_t$ = estimated network attack probability output by
the Intrusion Detection System (IDS) module.

\subsubsection{Action Space}

The action space mirrors the decision variable~\eqref{eq:decision_var}:
\begin{equation}
  \set{A} = \{0,\, 1,\, \ldots,\, M+2\},
  \quad |\set{A}| = M + 3.
  \label{eq:dqn_actions}
\end{equation}
Each action $a_t \in \set{A}$ selects an execution tier; the DQN outputs a
Q-value for each action, and the top-$K$ actions with highest Q-values form
the candidate subspace passed to BBO.

\subsubsection{CI-Modulated Reward Function}

The reward signal encodes the system objectives and penalises SLA
violations proportionally to patient criticality:
\begin{equation}
  r_t = -\!\left(
    \hat{w}_E \cdot \hat{E}_t
    + \hat{w}_L \cdot \hat{T}_t
    + \hat{w}_P \cdot \hat{R}_{P,t}
  \right)
  + \Phi_t \cdot \xi_{\mathrm{SLA}}(t),
  \label{eq:dqn_reward}
\end{equation}
where
\begin{equation}
  \xi_{\mathrm{SLA}}(t) =
  \begin{cases}
    -5 & \text{if } T_i(x_i) > T^{\max}_i, \\
     0 & \text{otherwise,}
  \end{cases}
  \label{eq:sla_penalty}
\end{equation}
and $\hat{w}_E, \hat{w}_L, \hat{w}_P$ are the normalised CI-adaptive weights
from~\eqref{eq:weight_normalised}.  Multiplying the SLA penalty by $\Phi_t$
ensures that violations carry maximum cost for critically ill patients
($\Phi_t \approx 1$) while being less severely penalised for routine tasks
($\Phi_t \approx 0$).

\subsubsection{Q-Value Update (Bellman Equation)}

The DQN parameters are updated online using the standard temporal-difference
target:
\begin{equation}
  Q(s_t, a_t) \leftarrow Q(s_t, a_t)
    + \eta \!\left[
        r_t + \gamma \max_{a'} Q(s_{t+1}, a') - Q(s_t, a_t)
      \right],
  \label{eq:bellman}
\end{equation}
where $\eta \in (0,1]$ is the learning rate and $\gamma \in [0,1)$ is the
discount factor governing the trade-off between immediate and future rewards.
In practice, the expectation in~\eqref{eq:bellman} is approximated via
mini-batch sampling from an experience replay buffer $\set{D}$ to
decorrelate training samples and improve convergence stability.

%% -----------------------------------------------------------------------
\subsection{BBO-DRL Hybridisation}
\label{subsec:hybrid}
%% -----------------------------------------------------------------------

The two components interact within each scheduling epoch as follows:
\begin{enumerate}
  \item \textbf{State observation}: the DQN agent reads $\vect{s}_t$
        from~\eqref{eq:dqn_state}.
  \item \textbf{Candidate pre-selection}: the DQN evaluates $Q(\vect{s}_t, a)$
        for all $a \in \set{A}$ and selects the top-$K$ candidate nodes
        $\set{A}_K \subseteq \set{A}$ with highest Q-values.
  \item \textbf{BBO subspace search}: BBO runs for $T_{\mathrm{BBO}}$
        iterations restricted to $\set{A}_K$, minimising $F(\cdot)$ from
        ~\eqref{eq:cost_function} subject to~\eqref{eq:constraint_sla}--\eqref{eq:constraint_domain},
        and returns $x_i^* = \argmin_{x \in \set{A}_K} F(x)$.
  \item \textbf{Execution}: task $\tau_i$ is dispatched to node $x_i^*$.
  \item \textbf{Policy update}: the environment transitions to
        $\vect{s}_{t+1}$; the tuple $(\vect{s}_t, x_i^*, r_t, \vect{s}_{t+1})$
        is stored in $\set{D}$; the DQN performs a gradient update
        via~\eqref{eq:bellman}.
\end{enumerate}

\subsubsection{Per-Task Computational Complexity}

Each scheduling epoch incurs a cost of $\mathcal{O}(|\set{A}|)$ for the DQN
forward pass, plus $\mathcal{O}(N_{\mathrm{pop}} \cdot T_{\mathrm{BBO}} \cdot K)$
for BBO search within $\set{A}_K$.  Since $K \ll M$ and $T_{\mathrm{BBO}}$
is small (empirically $\leq 20$ iterations), the overall per-task complexity is
\begin{equation}
  \mathcal{O}\!\left(N_{\mathrm{pop}} \cdot T_{\mathrm{BBO}} \cdot M\right),
  \label{eq:complexity}
\end{equation}
which scales linearly with the number of available compute nodes $M$,
making the scheduler suitable for real-time operation in clinical IoMT
deployments with $M \leq 50$ edge/fog nodes.

%% -----------------------------------------------------------------------
%% References (placeholder for standalone compilation)
%% -----------------------------------------------------------------------
\begin{thebibliography}{9}

\bibitem{burd1995processor}
T.~D. Burd and R.~W. Brodersen,
``Processor design for portable systems,''
\emph{J. VLSI Signal Process.}, vol.~13, no.~2--3, pp.~203--221, Aug. 1995.

\bibitem{shannon1948}
C.~E. Shannon,
``A mathematical theory of communication,''
\emph{Bell Syst. Tech. J.}, vol.~27, no.~3, pp.~379--423, Jul. 1948.

\bibitem{kleinrock1975}
L.~Kleinrock,
\emph{Queuing Systems, Volume 1: Theory}.
New York, NY, USA: Wiley-Interscience, 1975.

\bibitem{mnih2015}
V.~Mnih \emph{et al.},
``Human-level control through deep reinforcement learning,''
\emph{Nature}, vol.~518, no.~7540, pp.~529--533, Feb. 2015.

\bibitem{rappaport2002}
T.~S. Rappaport,
\emph{Wireless Communications: Principles and Practice}, 2nd~ed.
Upper Saddle River, NJ, USA: Prentice Hall, 2002.

\end{thebibliography}

\end{document}
